
\medskip
\phantomsection\label{par:sga3-VIB-11.4}%
\noindent\textbf{11.4.---}%
\footnote{\label{note:sga3-VIB-108}Editor's note: in
\SGARefLink{par:sga3-VIB-11.4}{11.4}--\SGARefLink{par:sga3-VIB-11.6}{11.6}, we have considered
\(\underline{\operatorname{Hom}}_{\mathcal O_S}
  (\mathrm W(\mathcal E),\mathrm W(\mathcal F))\)
instead of
\(\underline{\operatorname{Hom}}_{\mathcal O_S}
  (\mathrm V(\mathcal F),\mathrm V(\mathcal E))\)
\textup{(cf.\ I, \SGARefLink{sga3:i:proposition:4-6-3}{4.6.3})}, and simplified the original by
taking account of the addition \hyperref[concl:sga3-VIB-11.0-b]{11.0}\textup{(b)}.}
Let \(\mathcal E\) and \(\mathcal F\) be two
quasi-coherent \(\mathcal O_S\)-modules. Consider the
\(S\)-functor
\[
  \underline{\operatorname{Hom}}_{\mathcal O_S}
    \bigl(\mathrm W(\mathcal E),\mathrm W(\mathcal F)\bigr)
\]
\textup{(cf.\ I, \SGARefLink{sga3:i:definition:3-1-4}{3.1.4})}, that is, for every \(S\)-scheme
\(f:X\to S\),
\phantomsection\label{form:sga3-VIB-11.4-HomX}%
\[
  \underline{\operatorname{Hom}}_{\mathcal O_S}
    \bigl(\mathrm W(\mathcal E),\mathrm W(\mathcal F)\bigr)(X)
  =
  \operatorname{Hom}_{\mathcal O_X}
    \bigl(\mathrm W(\mathcal E)_X,\mathrm W(\mathcal F)_X\bigr).
\]
Moreover, by I, \SGARefLink{sga3:i:proposition:4-6-2}{4.6.2}, one has
\(\mathrm W(\mathcal E)_X=\mathrm W(f^*(\mathcal E))\)
\textup{(and similarly for \(\mathcal F\))}, and
\[
  \operatorname{Hom}_{\mathcal O_X}
    \bigl(\mathrm W(f^*(\mathcal E)),
          \mathrm W(f^*(\mathcal F))\bigr)
  =
  \operatorname{Hom}_{\mathcal O_X}
    \bigl(f^*(\mathcal E),f^*(\mathcal F)\bigr).
\]
We therefore obtain \textup{(using the adjunction formula
for the last equality)}
\phantomsection\label{form:sga3-VIB-11.4-adjunction}%
\[
  \tag{\(\dagger\)}
  \underline{\operatorname{Hom}}_{\mathcal O_S}
    \bigl(\mathrm W(\mathcal E),\mathrm W(\mathcal F)\bigr)(X)
  =
  \operatorname{Hom}_{\mathcal O_X}
    \bigl(f^*(\mathcal E),f^*(\mathcal F)\bigr)
  =
  \operatorname{Hom}_{\mathcal O_S}
    \bigl(\mathcal E,f_*f^*(\mathcal F)\bigr).
\]

\medskip
\phantomsection\label{prop:sga3-VIB-11.5}%
\noindent\textbf{Proposition 11.5.---}
\emph{Let \(X\) be an \(S\)-scheme that is quasi-compact and
quasi-separated over \(S\), let \(f:X\to S\) be the
structural morphism, and let \(\mathcal E\) and
\(\mathcal E'\) be two quasi-coherent
\(\mathcal O_S\)-modules. Suppose that one of the following
two conditions holds:}
\begin{enumerate}
\item[\textup{a)}]\phantomsection\label{concl:sga3-VIB-11.5-a}
  \emph{\(f\) is affine;}

\item[\textup{b)}]\phantomsection\label{concl:sga3-VIB-11.5-b}
  \emph{\(\mathcal E\) is flat over \(\mathcal O_S\).}
\end{enumerate}
\emph{Then the canonical morphism}
\phantomsection\label{form:sga3-VIB-11.5-projection}%
\[
  \mathcal E\otimes_{\mathcal O_S}\mathcal A(X)
    \longrightarrow f_*f^*(\mathcal E)
\]
\emph{is an isomorphism, and hence}
\phantomsection\label{form:sga3-VIB-11.5-Hom}%
\[
  \underline{\operatorname{Hom}}_{\mathcal O_S}
    \bigl(\mathrm W(\mathcal E'),\mathrm W(\mathcal E)\bigr)(X)
  =
  \operatorname{Hom}_{\mathcal O_S}
    \bigl(\mathcal E',
          \mathcal E\otimes_{\mathcal O_S}\mathcal A(X)\bigr).
\]

Indeed, the second assertion follows from
\hyperref[par:sga3-VIB-11.4]{11.4} and the first; the latter
follows from EGA~II, 1.4.7 in
case~\hyperref[concl:sga3-VIB-11.5-a]{\textup{(a)}}, and from
\hyperref[concl:sga3-VIB-11.0-b]{11.0\textup{(b)}} in
case~\hyperref[concl:sga3-VIB-11.5-b]{\textup{(b)}}.% 
\footnote{\label{note:sga3-VIB-109}Editor's note: we have
simplified the original, which used the isomorphism
\[
  \underline{\operatorname{Hom}}_{\mathcal O_S}
    \bigl(\mathrm W(\mathcal E'),\mathrm W(\mathcal E)\bigr)
  \simeq
  \underline{\operatorname{Hom}}_{\mathcal O_S}
    \bigl(\mathrm V(\mathcal E),\mathrm V(\mathcal E')\bigr),
\]
then the inclusion of the right-hand term in
\[
  \underline{\operatorname{Hom}}_S
    \bigl(\mathrm V(\mathcal E),\mathrm V(\mathcal E')\bigr)
  =
  \operatorname{Hom}_{\mathcal O_S}
    \bigl(\mathcal E',
          f_*f^*(\operatorname{Sym}(\mathcal E))\bigr),
\]
and applied EGA~III, 4.1.15 to
\(\mathrm V(\mathcal E)=
  \operatorname{Spec}(\operatorname{Sym}(\mathcal E))\)
to deduce \SGARefLink{par:sga3-VIB-11.0}{11.0}\textup{\SGARefLink{par:sga3-VIB-11.0:item:b}{(b)}}.}

\medskip
\phantomsection\label{par:sga3-VIB-11.6}%
\noindent\textbf{11.6.---}%
\footnote{\label{note:sga3-VIB-110}Editor's note: we have
expanded \SGARefLink{par:sga3-VIB-11.6}{11.6} and displayed the results obtained in the form
of \hyperref[prop:sga3-VIB-11.6.1]{Proposition 11.6.1}.}
Let \(G\) be an \(S\)-group and let \(\mathcal E\) be a
quasi-coherent \(\mathcal O_S\)-module. Giving \(\mathcal E\)
a \(G\)-\(\mathcal O_S\)-module structure \textup{(that is,
an \(\mathcal O_S\)-linear action of \(G\) on
\(\mathrm W(\mathcal E)\), cf.\ I, \SGARefLink{sga3:i:definition:4-7-1}{4.7.1})} is equivalent to
giving a morphism of \(S\)-functors in monoids
\phantomsection\label{form:sga3-VIB-11.6-rho}%
\[
  \rho:G\longrightarrow
  \underline{\operatorname{End}}_{\mathcal O_S}
    \bigl(\mathrm W(\mathcal E)\bigr);
\]
indeed, such a \(\rho\) necessarily takes \(G\) into
\(\underline{\operatorname{Aut}}_{\mathcal O_S}
  (\mathrm W(\mathcal E))\).

By \hyperref[par:sga3-VIB-11.4]{11.4}, giving a morphism of
\(S\)-functors
\[
  \rho:G\longrightarrow
  \underline{\operatorname{End}}_{\mathcal O_S}
    \bigl(\mathrm W(\mathcal E)\bigr)
\]
is equivalent to giving an element \(\theta\) of
\[
  \operatorname{Hom}_{\mathcal O_G}
    \bigl(f^*(\mathcal E),f^*(\mathcal E)\bigr),
\]
which corresponds by adjunction to an element \(\delta\) of
\[
  \operatorname{Hom}_{\mathcal O_S}
    \bigl(\mathcal E,f_*f^*(\mathcal E)\bigr),
\]
where \(f\) denotes the projection \(G\to S\).

Let \(m:G\times_SG\to G\) be multiplication, let
\(\delta_G\) be the morphism
\[
  \mathcal O_G\longrightarrow
  m_*(\mathcal O_{G\times_SG}),
\]
and let \(\varphi:G\times_SG\to S\) be the projection
\textup{(which equals \(f\circ m\))}. It is convenient to
write \(\boxtimes\) for the ``external'' tensor product. We
thus obtain a morphism
\phantomsection\label{form:sga3-VIB-11.6-id-deltaG}%
\[
  \operatorname{id}_{\mathcal E}\boxtimes\delta_G:
  f^*(\mathcal E)
    =\mathcal E\boxtimes_{\mathcal O_S}\mathcal O_G
  \longrightarrow
  \mathcal E\boxtimes_{\mathcal O_S}
    m_*(\mathcal O_{G\times_SG}).
\]
By abuse of notation, we again write
\(\operatorname{id}_{\mathcal E}\boxtimes\delta_G\) for the
composite of this morphism with the canonical morphism
\[
  \mathcal E\boxtimes_{\mathcal O_S}
    m_*(\mathcal O_{G\times_SG})
  \longrightarrow
  m_*m^*f^*(\mathcal E)
  =m_*\varphi^*(\mathcal E).
\]

On the other hand, let \(h:G\to S\) denote a second copy of
\(f:G\to S\), and consider the following commutative diagram,
in which \(p,q\) denote the two projections:
\begin{center}
  \phantomsection\label{diag:sga3-VIB-11.6-projections}%
  \begin{tikzcd}[row sep=large,column sep=large]
    G\times_SG
      \arrow[r,"q"]
      \arrow[d,"p"']
      \arrow[dr,"\varphi" description] &
    G \arrow[d,"h"] \\
    G \arrow[r,"f"'] &
    S
  \end{tikzcd}
\end{center}
Writing \(\delta\) again for the morphism
\(\mathcal E\to h_*h^*(\mathcal E)\), we obtain the morphism
\phantomsection\label{form:sga3-VIB-11.6-delta-id}%
\[
  \delta\boxtimes\operatorname{id}_{\mathcal O_G}:
  f^*(\mathcal E)
    =\mathcal E\boxtimes_{\mathcal O_S}\mathcal O_G
  \longrightarrow
  h_*h^*(\mathcal E)\boxtimes_{\mathcal O_S}\mathcal O_G
    =f^*h_*h^*(\mathcal E).
\]
By abuse of notation, we again write
\(\delta\boxtimes\operatorname{id}_{\mathcal O_G}\) for the
composite of this morphism with the canonical morphism
\[
  f^*h_*h^*(\mathcal E)\longrightarrow
  p_*\varphi^*(\mathcal E).
\]

The condition that \(\rho\) be compatible with multiplication
is then equivalent to saying that, for every open subset
\(U\) of \(S\), the following diagram is commutative:
\begin{center}
  \phantomsection\label{diag:sga3-VIB-11.6-1}%
  \makebox[0pt][r]{\textup{(1)}\qquad}%
  {\scriptsize
  \begin{tikzcd}[row sep=large,column sep=large]
    \Gamma(U,\mathcal E)
      \arrow[r,"\delta"]
      \arrow[d,"\delta"'] &
    \Gamma\!\left(
      U\times_SG,\,
      \mathcal E\boxtimes_{\mathcal O_S}\mathcal O_G
    \right)
      \arrow[d,"{\operatorname{id}_{\mathcal E}\boxtimes\delta_G}"] \\
    \Gamma\!\left(
      U\times_SG,\,
      \mathcal E\boxtimes_{\mathcal O_S}\mathcal O_G
    \right)
      \arrow[r,"{\delta\boxtimes\operatorname{id}_{\mathcal O_G}}"'] &
    \Gamma\!\left(
      U\times_SG\times_SG,\,
      \mathcal E\boxtimes_{\mathcal O_S}\mathcal O_G
        \boxtimes_{\mathcal O_S}\mathcal O_G
    \right)
  \end{tikzcd}}
\end{center}
Furthermore, the unit section \(\varepsilon:S\to G\) induces
a morphism \(u\) from \(\mathcal O_G\) to
\[
  \varepsilon_*\varepsilon^*(\mathcal O_G)
  =\varepsilon_*(\mathcal O_S),
\]
and the condition that \(\rho\) preserve unit elements is
equivalent to the commutativity of the diagram
\begin{center}
  \phantomsection\label{diag:sga3-VIB-11.6-2}%
  \makebox[0pt][r]{\textup{(2)}\qquad}%
  {\small
  \begin{tikzcd}[row sep=large,column sep=large]
    \Gamma(U,\mathcal E)
      \arrow[rr,"\delta"]
      \arrow[dr,"\sim"'] &
    {} &
    \Gamma\!\left(
      U\times_SG,\,
      \mathcal E\boxtimes_{\mathcal O_S}\mathcal O_G
    \right)
      \arrow[dl,"{\operatorname{id}_{\mathcal E}\boxtimes u}"] \\
    {} &
    \Gamma\!\left(
      U\times_SS,\,
      \mathcal E\boxtimes_{\mathcal O_S}\mathcal O_S
    \right)\mathpunct{.} &
    {}
  \end{tikzcd}}
\end{center}

Thus, giving \(\mathcal E\) a
\(G\)-\(\mathcal O_S\)-module structure is equivalent to
giving a morphism of \(\mathcal O_S\)-modules
\[
  \delta:\mathcal E\longrightarrow f_*f^*(\mathcal E)
\]
satisfying conditions
\hyperref[diag:sga3-VIB-11.6-1]{\textup{(1)}} and
\hyperref[diag:sga3-VIB-11.6-2]{\textup{(2)}} above. In that
case, the morphism
\[
  \theta:f^*(\mathcal E)\longrightarrow f^*(\mathcal E)
\]
deduced from \(\delta\) by adjunction is an isomorphism,
because it corresponds to the isomorphism
\[
  G\times_S\mathrm W(\mathcal E)
    \xrightarrow{\ \sim\ }
  G\times_S\mathrm W(\mathcal E)
\]
defined set-theoretically by
\((g,x)\longmapsto(g,gx)\); see also I, \SGARefLink{sga3:i:remark:6-5-4}{6.5.4}.

Now suppose that \(G\) is flat, quasi-compact, and
quasi-separated over \(S\), and that \(\mathcal A(G)\) is a
flat \(\mathcal O_S\)-module. Then, by
\hyperref[concl:sga3-VIB-11.1-c]{11.1\textup{(c)}}, the
canonical morphism
\[
  \mathcal A(G)\otimes_{\mathcal O_S}\mathcal A(G)
    \longrightarrow
  \mathcal A(G\times_SG)
\]
is an isomorphism, and we write \(\Delta\) for the morphism
\[
  \delta_G:\mathcal A(G)\longrightarrow
  \mathcal A(G)\otimes_{\mathcal O_S}\mathcal A(G).
\]

If, moreover,
\[
  \mathcal E\otimes_{\mathcal O_S}\mathcal A(G)
    =f_*f^*(\mathcal E),
\]
as is the case by \hyperref[prop:sga3-VIB-11.5]{11.5} if
\(G\to S\) is affine or if \(\mathcal E\) is flat over
\(\mathcal O_S\), then conditions
\hyperref[diag:sga3-VIB-11.6-1]{\textup{(1)}} and
\hyperref[diag:sga3-VIB-11.6-2]{\textup{(2)}} are equivalent
to the following conditions, which express that
\[
  \delta:\mathcal E\longrightarrow
  \mathcal E\otimes_{\mathcal O_S}\mathcal A(G)
\]
makes \(\mathcal E\) a right \(\mathcal A(G)\)-comodule
\textup{(cf.\ I, \SGARefLink{sga3:i:proposition:4-7-2}{4.7.2})}:

\smallskip
\phantomsection\label{cond:sga3-VIB-11.6-CM1}%
\noindent\textup{(CM 1)}
Put \(\mathcal A=\mathcal A(G)\). The following diagram is
commutative:
\begin{center}
  \phantomsection\label{diag:sga3-VIB-11.6-CM1}%
  \begin{tikzcd}[row sep=large,column sep=large]
    \mathcal E
      \arrow[r,"\delta"]
      \arrow[d,"\delta"'] &
    \mathcal E\otimes\mathcal A
      \arrow[d,"{\operatorname{id}_{\mathcal E}\otimes\Delta}"] \\
    \mathcal E\otimes\mathcal A
      \arrow[r,"{\delta\otimes\operatorname{id}_{\mathcal A}}"'] &
    \mathcal E\otimes\mathcal A\otimes\mathcal A
  \end{tikzcd}
\end{center}

\smallskip
\phantomsection\label{cond:sga3-VIB-11.6-CM2}%
\noindent\textup{(CM 2)}
Write
\(\eta:\mathcal A\to\mathcal O_S\) for the morphism
\(\mathcal A(\varepsilon)\). The following diagram is
commutative:
\begin{center}
  \phantomsection\label{diag:sga3-VIB-11.6-CM2}%
  \begin{tikzcd}[row sep=large,column sep=large]
    \mathcal E
      \arrow[rr,"\delta"]
      \arrow[dr,"\sim"'] &
    {} &
    \mathcal E\otimes\mathcal A
      \arrow[dl,"{\operatorname{id}_{\mathcal E}\otimes\eta}"] \\
    {} &
    \mathcal E\otimes\mathcal O_S\mathpunct{.} &
    {}
  \end{tikzcd}
\end{center}

\medskip
\phantomsection\label{rem:sga3-VIB-11.6.A}%
\noindent\textbf{Remark 11.6.A.---}
Recall that \(\mathrm V(\mathcal E)\) denotes the vector
fibration over \(S\) representing the functor
\(\mathrm V(\mathcal E)\); that is, for every \(S'\to S\),
\[
  \mathrm V(\mathcal E)(S')
  =
  \operatorname{Hom}_{\mathcal O_{S'}}
    \bigl(\mathcal E_{S'},\mathcal O_{S'}\bigr).
\]
Since I, \SGARefLink{sga3:i:proposition:4-6-2}{4.6.2} gives an anti-isomorphism of \(S\)-functors
in monoids
\[
  \underline{\operatorname{End}}_{\mathcal O_S}
    \bigl(\mathrm W(\mathcal E)\bigr)
  \simeq
  \underline{\operatorname{End}}_{\mathcal O_S}
    \bigl(\mathrm V(\mathcal E)\bigr),
\]
if \(\mathcal E\) is a left
\(G\)-\(\mathcal O_S\)-module, then there is a right action
\[
  \mu:\mathrm V(\mathcal E)\times_SG
    \longrightarrow\mathrm V(\mathcal E)
\]
of \(G\) on \(\mathrm V(\mathcal E)\), defined
set-theoretically by
\[
  (\phi g)(x)=\phi(gx)
\]
for every \(g\in G(S')\),
\(x\in\Gamma(S',\mathcal E_{S'})\), and
\(\phi\in
  \operatorname{Hom}_{\mathcal O_{S'}}
    (\mathcal E_{S'},\mathcal O_{S'})\).
We therefore have the commutative diagrams
\begin{center}
  \phantomsection\label{diag:sga3-VIB-11.6.A-actions}%
  {\small
  \begin{tabular}{c@{\qquad}c}
    \begin{tikzcd}[row sep=large,column sep=large]
      \mathrm V(\mathcal E)\times_SG\times_SG
        \arrow[r,"{\mu\times\operatorname{id}_G}"]
        \arrow[d,"{\operatorname{id}_{\mathrm V(\mathcal E)}\times m}"'] &
      \mathrm V(\mathcal E)\times_SG
        \arrow[d,"\mu"] \\
      \mathrm V(\mathcal E)\times_SG
        \arrow[r,"\mu"'] &
      \mathrm V(\mathcal E)
    \end{tikzcd}
    &
    \begin{tikzcd}[row sep=large,column sep=large]
      \mathrm V(\mathcal E) &
      \mathrm V(\mathcal E)\times_SG
        \arrow[l,"\mu"'] \\
      {} &
      \mathrm V(\mathcal E)\times_SS\mathpunct{.}
        \arrow[u,"{\operatorname{id}_{\mathrm V(\mathcal E)}\times\varepsilon}"]
        \arrow[ul,"\sim"']
    \end{tikzcd}
  \end{tabular}}
\end{center}
When \(G\) is flat, quasi-compact, and quasi-separated over
\(S\), when \(\mathcal A(G)\) is a flat
\(\mathcal O_S\)-module, and when one of the conditions of
\hyperref[prop:sga3-VIB-11.5]{11.5} is satisfied, one
similarly recovers conditions
\hyperref[cond:sga3-VIB-11.6-CM1]{\textup{(CM 1)}} and
\hyperref[cond:sga3-VIB-11.6-CM2]{\textup{(CM 2)}}.

Consequently, we have proved:

\medskip
\phantomsection\label{prop:sga3-VIB-11.6.1}%
\noindent\textbf{Proposition 11.6.1.---}
\emph{Let \(G\) be an \(S\)-group that is flat,
quasi-compact, and quasi-separated over \(S\), such that
\(\mathcal A(G)\) is a flat \(\mathcal O_S\)-module, and let
\(\mathcal E\) be a quasi-coherent
\(\mathcal O_S\)-module.}
\begin{enumerate}
\item[\textup{(i)}]\phantomsection\label{concl:sga3-VIB-11.6.1-i}
  \emph{Giving an \(\mathcal A(G)\)-comodule structure}
  \[
    \delta:\mathcal E\longrightarrow
      \mathcal E\otimes_{\mathcal O_S}\mathcal A(G)
  \]
  \emph{is equivalent to giving
  \(\mathcal E\) a
  \(G_{\mathrm{af}}\)-\(\mathcal O_S\)-module structure
  \textup{(that is, an \(\mathcal O_S\)-linear action of
  \(G_{\mathrm{af}}\) on \(\mathcal E\))}. By composition
  with the morphism of \(S\)-groups
  \(G\to G_{\mathrm{af}}\), this defines a
  \(G\)-\(\mathcal O_S\)-module structure on
  \(\mathcal E\).}

\item[\textup{(ii)}]\phantomsection\label{concl:sga3-VIB-11.6.1-ii}
  \emph{If, moreover, \(\mathcal E\) is flat, every
  \(\mathcal O_S\)-linear action of \(G\) on
  \(\mathcal E\) factors through \(G_{\mathrm{af}}\) and
  corresponds to a unique \(\mathcal A(G)\)-comodule
  structure on \(\mathcal E\).}
\end{enumerate}

\medskip
\phantomsection\label{lem:sga3-VIB-11.7}%
\noindent\textbf{Lemma 11.7.---}
\emph{Let \(G\) be an \(S\)-group that is flat,
quasi-compact, and quasi-separated over \(S\), such that
\(\mathcal A=\mathcal A(G)\) is a flat
\(\mathcal O_S\)-module. Let \(\mathcal E\) be a
quasi-coherent \(\mathcal O_S\)-module, let}
\[
  \delta:\mathcal E\longrightarrow\mathcal E\otimes\mathcal A
\]
\emph{be an \(\mathcal A\)-comodule structure, and let}
\[
  \rho:G\longrightarrow
  \underline{\operatorname{Aut}}_{\mathcal O_S}
    \bigl(\mathrm W(\mathcal E)\bigr)
\]
\emph{be the associated action of \(G\) on \(\mathcal E\).
Let \(\mathcal E_0\) be a quasi-coherent
\(\mathcal O_S\)-submodule of \(\mathcal E\) such that the
restriction \(\delta_0\) of \(\delta\) to
\(\mathcal E_0\)}%
\begingroup
\renewcommand{\thefootnote}{\fnsymbol{footnote}}%
\footnote[1]{\label{note:sga3-VIB-11.7-delta-source}
Translator's note: the French source writes ``the restriction
\(\delta_0\) of \(\mu\)'' here. The surrounding notation and
the displayed square show that \(\delta\) is intended.}%
\endgroup
\emph{ factors through}
\[
  \mathcal E_0\longrightarrow\mathcal E_0\otimes\mathcal A;
\]
\emph{that is, suppose that the following diagram is
commutative:}
\begin{center}
  \phantomsection\label{diag:sga3-VIB-11.7-stability}%
  \begin{tikzcd}[row sep=large,column sep=large]
    \mathcal E_0
      \arrow[r,hook]
      \arrow[d,"\delta_0"'] &
    \mathcal E
      \arrow[d,"\delta"] \\
    \mathcal E_0\otimes\mathcal A
      \arrow[r,hook] &
    \mathcal E\otimes\mathcal A\mathpunct{.}
  \end{tikzcd}
\end{center}
\emph{\textup{(N.B.\ The morphism
\(\mathcal E_0\otimes\mathcal A
  \to\mathcal E\otimes\mathcal A\)
is injective because \(\mathcal A\) is flat over
\(\mathcal O_S\).)}
Then \(\delta_0\) makes \(\mathcal E_0\) an
\(\mathcal A(G)\)-comodule, and hence defines an action
\(\rho_0\) of \(G\) on \(\mathcal E_0\). This is called the
action on \(\mathcal E_0\) induced by \(\rho\), and we say
that \(\mathcal E_0\) is stable under \(\rho\).}

This follows immediately from the definitions and from
\hyperref[par:sga3-VIB-11.6]{11.6}. Notice, however, that in
general the canonical map
\[
  \mathrm W(\mathcal E_0)\longrightarrow\mathrm W(\mathcal E)
\]
is not a monomorphism.

\medskip
\phantomsection\label{rem:sga3-VIB-11.7.bis}%
\noindent\textbf{Remark 11.7.bis.---}%
\footnote{\label{note:sga3-VIB-111}Editor's note: we have
added this remark, which generalizes \SGARefLink{lem:sga3-VIB-11.7}{11.7} and will be useful
in \SGARefLink{prop:sga3-VIB-11.10.bis}{11.10.b}is.}
Let \(G\) be a flat \(S\)-group and let \(\mathcal E\) be a
quasi-coherent \(G\)-\(\mathcal O_S\)-module. Write \(f\) for
the morphism \(G\to S\) and \(\delta\) for the morphism
\[
  \mathcal E\longrightarrow f_*f^*(\mathcal E)
\]
defined in \hyperref[par:sga3-VIB-11.6]{11.6}. Let
\(\mathcal E_0\) be a quasi-coherent
\(\mathcal O_S\)-submodule of \(\mathcal E\). Since \(f\) is
flat, \(f_*f^*(\mathcal E_0)\) is an
\(\mathcal O_S\)-submodule of \(f_*f^*(\mathcal E)\), and the
same holds for \(\varphi=f\times f\). Consequently, if the
restriction \(\delta_0\) of \(\delta\) to \(\mathcal E_0\)
factors through \(f_*f^*(\mathcal E_0)\), then it makes
\(\mathcal E_0\) a \(G\)-\(\mathcal O_S\)-module. In this
case, we say that \(\mathcal E_0\) is a \(G\)-stable
submodule of \(\mathcal E\).
